% ============================================================
%  CSE200 Online-1 (C2)  ->  "CSE200: Online-1 on LaTeX"
%  Compile:  pdflatex Q2_LaTeXIntro.tex
%  Needs the `mwe` package for the example-image-a/b/c demo images.
% ============================================================
\documentclass[11pt]{article}

\usepackage[T1]{fontenc}
\usepackage[utf8]{inputenc}
\usepackage[margin=1in]{geometry}
\usepackage{amsmath}
\usepackage{amssymb}
\usepackage{graphicx}
\usepackage{subcaption}
\usepackage{multirow}
\usepackage{mwe}

\title{CSE200: Online-1 on \LaTeX}
\author{Your Name (Student ID)}
\date{January 3, 2026}

\begin{document}
\maketitle

% ------------------------------------------------------------
\section{Introduction}
Document preparation systems are essential for producing structured and
professional technical documents. \textbf{\LaTeX} allows authors to focus on
content while maintaining consistent formatting. This article demonstrates
formatted text, lists, equations, tables, and figures using commands taught in
class.

\subsection{Text Emphasis and Fonts}
This sentence contains \textbf{bold text}, \textit{italic text},
\emph{emphasized text}, and \underline{underlined text}. Font sizes can also vary
such as {\Large Large text}, {\small small text}, and {\huge huge text}.

% ------------------------------------------------------------
\section{List Structures}

\subsection{Mixed and Nested Lists}
\begin{itemize}
    \item Primary Feature
    \item Secondary Features
    \begin{enumerate}
        \item Formatting control
        \item Mathematical typesetting
        \begin{itemize}
            \item Inline math
            \item Display math
        \end{itemize}
    \end{enumerate}
    \item[\textbf{Note}] Lists can be customized.
\end{itemize}

% ------------------------------------------------------------
\section{Mathematical Modeling}
Mathematical expressions are often used to represent data relationships. Consider
variables $x_i$, $y$, and $z^2$.

% ------------------------------------------------------------
\section{Mathematical Modeling}
Mathematical expressions are commonly used to describe relationships among
variables. Let the variables be $x_i$, $y_j$, and $\sigma^2$.

\subsection{Displayed Equations}
\begin{equation}
    y_j = \alpha x_j^2 + \beta x_j + \gamma
\end{equation}
\begin{equation}
    \begin{split}
        T &= \sum_{i=1}^{n}(x_i-\mu)^2 \\
          &= \sum_{i=1}^{n} x_i^2 - 2\mu\sum_{i=1}^{n} x_i + n\mu^2
    \end{split}
\end{equation}
\[
    f(x_i) = \begin{cases}
        \dfrac{x_i^2}{\sigma^2} & \text{if } x_i \ge 0 \\[2ex]
        \dfrac{|x_i|}{\sigma}   & \text{if } x_i < 0
    \end{cases}
\]

% ------------------------------------------------------------
\section{Tabular Data Presentation}
\begin{table}[ht]
    \centering
    \begin{tabular}{|c|c|c|}
        \hline
        \multirow{2}{*}{\textbf{Module}} & \multicolumn{2}{c|}{\textbf{Score}} \\
        \cline{2-3}
                 & Theory & Lab \\
        \hline
        Module A & 75 & 80 \\
        \hline
        \multirow{2}{*}{Module B} & 88 & 90 \\
                                  & 85 & 87 \\
        \hline
        Module C & 70 & 72 \\
        \hline
    \end{tabular}
    \caption{Module-wise Performance Summary}
    \label{tab:modules}
\end{table}

% ------------------------------------------------------------
\section{Visual Comparison}

\subsection{Asymmetric Subfigure Layout}
\begin{figure}[ht]
    \centering
    \begin{subfigure}{0.6\textwidth}
        \centering
        \includegraphics[width=\textwidth]{example-image-a}
        \caption{Main Diagram}
    \end{subfigure}

    \vspace{1em}
    \begin{subfigure}{0.35\textwidth}
        \centering
        \includegraphics[width=\textwidth]{example-image-b}
        \caption{Detail B}
    \end{subfigure}
    \hfill
    \begin{subfigure}{0.35\textwidth}
        \centering
        \includegraphics[width=\textwidth]{example-image-c}
        \caption{Detail C}
    \end{subfigure}
    \caption{Asymmetric layout with one dominant figure}
    \label{fig:asym}
\end{figure}

Figure~\ref{fig:asym} illustrates how visual elements can be arranged
systematically within a document.

% ------------------------------------------------------------
\section*{Conclusion}
This article highlighted the importance of \textbf{structured writing} using
\LaTeX. The use of \textit{formatted text}, mathematical expressions, tables, and
figures improves both \underline{readability} and presentation quality.

\end{document}
